Para identificar las dificultades conceptuales que enfrentan los estudiantes de la asignatura FLP, se realizó un diagnóstico basado en tres fuentes complementarias: una encuesta dirigida a estudiantes que han cursado FLP en distintos períodos académicos, una encuesta aplicada a docentes del programa de Ingeniería de Sistemas y una revisión del microcurrículo de la asignatura.

Estas fuentes permitieron obtener una caracterización de las dificultades desde diferentes perspectivas, teniendo la de los estudiantes que las experimentan directamente, la de los docentes que se encargan de enseñar el contenido y la del diseño curricular que estructura el aprendizaje, lo que generó una selección justificada de los conceptos que servirán como base para la biblioteca de casos predefinidos del prototipo.

\section{Metodología}

El diagnóstico se apoyó en un enfoque mixto, combinando métodos cuantitativos provenientes de encuestas con escala Likert y métodos cualitativos a través de preguntas abiertas y análisis documental del microcurrículo.

\subsection{Diseño de los instrumentos}

El diseño de ambas encuestas se fundamentó en las dimensiones de evaluación identificadas en la literatura acerca de las dificultades en la enseñanza de compiladores e intérpretes \cite{Stamenkovic2024AWE} \cite{jovanovic2021comvis} y sobrecarga cognitiva en programación \cite{zagami2023}, que orientaron la selección de los ítems a evaluar: comprensión teórica, capacidad de aplicación, percepción de las herramientas, disponibilidad de recursos y valoración de alternativas interactivas. Por otro lado, el microcurrículo oficial de la asignatura permitió anclar las preguntas a los resultados de aprendizaje e indicadores de logro concretos de FLP, evitando que el instrumento midiera percepciones genéricas desvinculadas del contexto específico de la asignatura. 

La encuesta a estudiantes fue diseñada con siete ítems de escala Likert de cinco niveles (1: Totalmente en desacuerdo, 5: Totalmente de acuerdo) para evaluar la percepción de las dificultades conceptuales en FLP, junto con una pregunta abierta orientada a identificar las recomendaciones de mejora. Esta estructura permitió distinguir entre dificultades atribuibles al contenido, al entorno de programación y a la escasez de recursos complementarios. 

El diseño de la encuesta a docentes incluyó una pregunta de segmentación, teniendo en cuenta el semestre en el que cada docente dicta programación principalmente, seguida de bloques temáticos con ítems Likert sobre la percepción de las dificultades estudiantiles en distintos paradigmas de programación, cerrando con una pregunta abierta para comentarios o sugerencias respecto al proceso de enseñanza de programación. 

Finalmente, la revisión del microcurrículo se realizó sobre el documento oficial de la asignatura FLP (código 750017C) el cual se encuentra en el Anexo \ref{annex:microcurriculum}. Este documento estructura los contenidos a través de competencias específicas, resultados de aprendizaje (RA) e indicadores de logro (IL), lo que permitió contrastar las dificultades reportadas por estudiantes y docentes con la estructura formal de los contenidos e identificar cuáles indicadores concentran mayor complejidad conceptual. 

\subsection{Selección de la población}

La población objetivo de la encuesta estudiantil fueron estudiantes que hubieran cursado FLP en la Universidad del Valle sede Tuluá durante los últimos tres años académicos, optando por no restringir la convocatoria al período en curso, con el objetivo de ampliar el tamaño de la muestra y capturar perspectivas recientes y en retrospectiva, lo que permitió que el diagnóstico reflejara una experiencia acumulada con la asignatura. La encuesta fue aplicada de forma voluntaria, obteniendo 36 respuestas adjuntas en el Anexo \ref{annex:students}.

La población objetivo de la encuesta docente fueron los profesores del programa de Ingeniería de Sistemas que dictan asignaturas de programación, independientemente del semestre. Esta decisión busca obtener una perspectiva longitudinal del programa, considerando que mientras los docentes de semestres avanzados informan sobre las dificultades que observan directamente en FLP y asignaturas afines, los docentes de primeros semestres permiten caracterizar las bases con las que los estudiantes ingresan. La encuesta fue aplicada de forma voluntaria, obteniendo 5 respuestas, las cuales se encuentran detalladas en el Anexo \ref{annex:teachers}.